\subsection{Sharing sums across different common points}
\label{sec:cross-point-sharing}
Retain $h=46$, $F=\mathbb Q^h$ with form $I-J/9$, $v=15180$,
$N=v^3$, and $m=97336$. The following extension of the preceding computation
allows equal sums from different common-point groups to be identified.

\begin{proposition}[Complement frames for the reversed computation]
Consider a cancellation-free binary-addition graph, embedded reversibly by
reusing a pivot along one outgoing edge at each addition. Suppose the source
indicator span $U_z$ at every node $z$ is nondegenerate, and every source
contributing to a designated output is orthogonal to that output's target
indicator. Then the side computation admits increasing nondegenerate frames
in both stage orientations, with no extra decreasing edges.
\end{proposition}
\begin{proof}
In the forward middle computation, use
$D_{U_z}=(B\otimes F)\mathbin\perp(P\otimes U_z)$, with the future
nondegenerate line factor suppressed. Along every directed edge $z\to w$,
$U_z\subset U_w$. Physical roles follow paths, so these labels increase.
At an output, $U_z\subset t_T^\perp$, as required by the target injection.

In the reversed middle computation, use $D_{U_z^\perp}$ instead. The ambient
form is nondegenerate, so the orthogonal complement of every nondegenerate
$U_z$ is nondegenerate. Reversing an edge gives
$U_w^\perp\subset U_z^\perp$. At an original output, the physical $X$
line satisfies $\langle t_T\rangle\subset U_z^\perp$; at an original
input, the final complement is exactly the physical $Y$ source complement
$t_S^\perp$. These are the required reversed-stage endpoints.
All roles incident to a node gate receive the same frame. Newly active roles
enter from $D_0=B\otimes F$ and retired roles grow to $D_1=A\otimes F$.
Early mixers use $D_0$ and late cleanup mixers use $D_1$, as before.
Tensoring with the future line preserves these conclusions. Thus all side
edges are increasing in both orientations.
\end{proof}

This argument does not require the reachable-target spans to be positive
definite. It also does not permit arbitrary equal-value circuit rewrites:
nondegeneracy of the source spans and the stated endpoint orthogonality remain
necessary hypotheses of this transfer.

Start with the preceding $h$ copies of the $n=45$ pair-exclusion circuit.
Identify their input nodes when they name the same physical triple. Identify
addition nodes whenever their formal source-support sets agree, retaining
the first encountered decomposition. Process common points in increasing
order and the local node identifiers in increasing order, then discard unused
nodes. Keep all $q=h\binom{h-1}{2}=45540$ designated partial outputs. Their
injection into the physical target bank adds the three contributions for each
target triple. A neighboring pair has a unique common point, so the resulting
matrix remains precisely the intersection-one relation.

There is a compact exact method for finding every cross-group equality.
If a sum belongs to groups $i\ne j$, every source triple contains both $i$
and $j$. Its support is therefore a pair-star
$\{\{i,j,k\}:k\in K\}$, or a single input triple. The fixed pair and the
union of its vertices uniquely determine this support. All other addition
nodes belong to a single group and were already interned in the local circuit.
The retained global graph has
\[
 c=486330,\qquad q=45540,\qquad R=c+q=531870.
\]
It merges 45706 of the previous 532036 addition nodes. Each retained node has
exactly the same source support as one of the original common-point nodes.
Every such span is positive definite: if all its triples contain $i$, then
$\sum_j u_j=3u_i$ and
$\langle u,u\rangle=\sum_{j\ne i}u_j^2>0$ for $u\ne0$.
The preceding proposition therefore applies, including after cross-group
sharing. The physical reversible roles are compiled afresh from this graph;
they are not obtained by identifying simultaneously occupied scratch slots.

For completeness, if a graph has $c$ additions and $q$ output uses, it has
$2c+q$ outgoing uses in total. Input gates allocate one role per use, and each
addition identifies its pivot input with one outgoing role. The resulting
invertible gather-and-fanout product $L$ uses $c+q$ roles. With $V$ the source
copy, $J$ the output injection, and $G,R$ the retained central operations, use
\[
 L,\ J,\ L^{-1},\ R,\ V,\ G,\ R,\ L,\ J,\ L^{-1},\ G,\ V.
\]
The two side injections sum to $JLz+JL(z+Vx)=JLVx$, and the final cleanup
restores arbitrary scratch $z$. Centers are also restored. Since $JLV+RG=I$,
each invocation is the original bank shear. The reversed invocation uses the
reversed inverse schedule and the complement frames just proved.

The only decreasing transitions remain the original central returns, of total
dimension $L_{\rm dec}=3v^2h^2$. Every auxiliary role starts at $D_0$ and ends
at $D_1$, so the previous bijective first/third-stage sharing, including centers,
still applies without changing its joining labels. Thus
\begin{align*}
 W_{\rm b}^{\rm shared}&=2N+2v^2(R+h)=252137288620800,\\
 s_{\rm b}^{\rm shared}&=W_{\rm b}^{\rm shared}m-N+6v^2h^2
                         =24542034552800107200,\\
 W_{\rm b}^{\rm shared}m-s_{\rm b}^{\rm shared}&=572394081600.
\end{align*}
\begin{proposition}[Cross-group shared-computation bit interface]
\label{prop:shared-point-bit-interface}
The shared-point bit network exchanges the scalar banks, restores arbitrary
auxiliary inputs, and satisfies the original all-role endpoint identity with
$W_{\rm b}^{\rm shared}$ roles and total edge rank $s_{\rm b}^{\rm shared}$.
\end{proposition}
\begin{proof}
The scalar and frame arguments above preserve the original stage boundaries.
Forward, inverse, forward shears exchange the banks. All surviving auxiliary
roles have endpoint matrices zero and $I_m$ and are fixed by the exchange.
Telescoping gives $W_{\rm b}^{\rm shared}m-2N+2L_{\rm dec}$ total absolute
label change, and the unchanged rational source correction adds $N$.
The original common-frame and rational-shear transfer applies; all graphs,
bases, and gates are fixed finite data.
\end{proof}

\subsection{Explicit rational bounds for the shared-point recurrences}
Put $a_{\rm b}=203/10^{11}$, $a_{\rm c}=9/500000000000$, and $L_0=5743/500$.
The exact comparisons are
\[
 \eta_{\rm b}^{\rm shared}=\frac{27}{1157655136}>a_{\rm b}L_0,
 \qquad \eta_{\rm c}=\frac7{22253827054}>a_{\rm c}L_0.
\]
Together with the retained logarithm enclosure $\log m<L_0$ and
$e^{-x}>1-x$ for $x>0$, these imply the required strict recurrence bounds with
\begin{equation}\label{eq:explicit-motif-exponents}
 \tau=1-\frac{203}{10^{11}},\qquad \sigma=1-\frac9{500000000000}.
\end{equation}
